%!TeX root=../quali.tex

\palavraschave{Computação de alto desempenho, Impacto ambiental, Sustentabilidade, Pegada de carbono, Pegada hídrica, Escalonamento de tarefas, Avaliação de ciclo de vida, Suficiência}

\keywords{High-performance computing,Environmental impact,Sustainability,Carbon footprint,Water footprint,Job scheduling,Lifecycle assessment,Sufficiency}

\resumo{
A computação de alto desempenho impõe custos ambientais em carbono, água e materiais ao longo de todo o seu ciclo de vida, mas a área tradicionalmente equipara esses custos ao consumo operacional de energia. Este trabalho procura entender o que a área mede, o que ela deixa de fora e como sistemas de alto desempenho podem ser operados e avaliados quando as dimensões menos estudadas são incorporadas. Ele começa com um estudo de mapeamento sistemático de 62 relatos que cobrem dezesseis anos de pesquisa, classificados segundo impactos ambientais, estágios do ciclo de vida, locais do sistema e mecanismos de gestão e intervenção. A literatura se concentra em um único corredor dessa classificação: carbono, estudado no estágio operacional, no nível da infraestrutura e por meio de medição. Carbono aparece em 54 dos 62 relatos, materiais e resíduos em 12, uso de água em 5 e biodiversidade em 1, e nenhum relato trata mais de dois impactos ao mesmo tempo. O corredor acompanha a disponibilidade de dados mais do que o formato do problema, e são muitas as células que ele deixa vazias. O restante do trabalho segue duas delas. A primeira permanece na operação, onde escalonadores, plataformas e decisões de substituição de hardware raramente precificam o custo ambiental. Estendemos o SimGrid e o Batsim com um modelo de pegada por host que contabiliza carbono operacional, água local e remota e impactos incorporados sob sinais de rede elétrica variáveis no tempo, e o usamos para avaliar o Greenfilling, uma regra que suprime inícios por backfilling quando a intensidade da rede excede sua média móvel. Em 288 comparações pareadas sobre três redes nacionais e duas máquinas de produção, nenhuma mediana de cenário alterou a pegada alvo em 0,1\%, enquanto o slowdown limitado médio aumentou em todos os pares, de modo que uma política precisa precificar tanto a espera quanto a ação, e não apenas recusar. Uma análise offline das mesmas cargas de trabalho situa o teto do deslocamento de tarefas em 6,76\% do carbono operacional e 4,89\% da água, com deslocamento máximo de 24 horas. A segunda acompanha como os sistemas são julgados. O Green500 os ordena por desempenho por watt, uma razão que pode melhorar enquanto o impacto absoluto por trás dela cresce. O plano desenvolve mecanismos que precificam o custo ambiental no escalonador, nos estados de energia da plataforma, no momento da substituição de hardware e junto aos usuários, e uma avaliação orientada à suficiência que usa o Green500 como ponto de referência para expor o que uma razão por watt esconde.
}

\abstract{
High-performance computing carries environmental costs in carbon, water, and materials across its whole lifecycle, yet the field has traditionally equated those costs with operational energy consumption. This work tries to understand what the field measures, what it leaves out, and how HPC systems can be operated and evaluated once the less studied dimensions are brought in. It begins with a systematic mapping study of 62 reports covering sixteen years of research, classified along environmental impacts, lifecycle stages, system loci, and management and intervention levers. The literature concentrates in a single corridor of that classification: carbon, studied at the operational stage, at the facility level, and through measurement. Carbon appears in 54 of the 62 reports, materials and waste in 12, water use in 5, and biodiversity in 1, and no report addresses more than two impacts at once. The corridor traces the availability of data more than the shape of the problem, and the cells it leaves empty are many. The rest of the work follows two of them. The first stays with operations, where schedulers, platforms, and hardware-replacement decisions rarely price environmental cost. We extended SimGrid and Batsim with a per-host footprint model that accounts for operational carbon, on-site and off-site water, and embodied impacts under time-varying grid signals, and used it to evaluate Greenfilling, a rule that suppresses backfill starts when grid intensity exceeds its moving average. Across 288 paired comparisons over three national grids and two production machines, no scenario median moved the targeted footprint by 0.1\%, while mean bounded slowdown rose in every pair, so a policy has to price both waiting and acting rather than only refusing. An offline analysis of the same workloads places the ceiling of job shifting at 6.76\% of operational carbon and 4.89\% of water, at a maximum displacement of 24 hours. The second follows how systems are judged. The Green500 ranks them by performance per watt, a ratio that can improve while the absolute burden behind it grows. The plan develops mechanisms that price environmental cost in the scheduler, in the platform's power states, in replacement timing, and with the users, and a sufficiency-aware evaluation that uses the Green500 as a reference point to expose what a per-watt ratio hides.
}
